\section{Introduction}

Reading is the process by which we extract meaning from print. For skilled readers this feels instantaneous, but most accounts agree that word recognition draws on representations at multiple grains and levels of abstraction, from sublexical orthographic form to word-level meaning. Accounts disagree sharply, however, on whether these levels are engaged as a sequence of discrete stages or as a single, continuous process, and on whether an intermediate lexical representation exists at all. English morphology offers a useful case for studying how word recognition unfolds and the representations involved. This is because English spelling exhibits irregular and context-sensitive grapheme–phoneme mappings (e.g., 'c' pronounced /k/ in cat but /s/ in city). At the same time, English preserves regularities at higher representational levels, particularly in its morphological structure; words that share a stem or affix often exhibit consistent orthographic patterns and systematic semantic relationships (e.g., \textit{heal--health} share a stem; \textit{teacher}, \textit{baker}, and \textit{driver} share the agentive suffix \textit{-er}). As a result, beginning readers must rely on implicit learning to detect the probabilistic positional and contextual regularities that define the complex English orthography-to-phonology mapping.

Readers vary in how well they have learned these complex mappings, and in their sensitivity to orthographic structure \citep{perfettiReadingAbilityLexical2007, castlesHowDoesOrthographic2006, stanovichExposurePrintOrthographic1989}. This variability provides a means of identifying the type of information involved in a given lexical manipulation: if the effect of a manipulation varies systematically with a reader's orthographic or semantic profile, that relationship is informative about the level of representation involved. The present study extends this logic to morphological processing. We examine two dimensions of reading skill, Orthographic Sensitivity (OS) and Language Proficiency (LP), in relation to two lexical variables that have been argued to index distinct levels of morphological representation: base frequency and morphological family size. By relating these individual differences to the N250 and N400—ERP components argued to reflect temporally distinct phases of processing, from sublexical form analysis to semantic integration—we ask whether Orthographic Sensitivity and Language Proficiency differentially modulate form-based and meaning-based aspects of morphological processing.

Before turning to the study's specific aims, we address several issues that provide the necessary theoretical and empirical foundation. First, we consider what base frequency and family size are each argued to index, and why competing theoretical accounts converge in predicting that their effects should emerge at different points in time. Second, we consider how Orthographic Sensitivity and Language Proficiency emerge as dissociable dimensions of individual variation. Third, we examine why the N250 and N400 provide appropriate neural indices of the two processing phases in question.

\subsection{Base Frequency and Morphological Family Size}

Base frequency, the cumulative token frequency of a base morpheme across all its surface forms, and morphological family size, the number of distinct words sharing that base, both facilitate visual word recognition, and their effects are independent when appropriately controlled \citep{fordDerivationalMorphologyBase2010}. The two variables have been argued to index different levels of morphological representation. Base frequency affects the cumulative ease with which a base morpheme's orthographic form is processed across its surface variants, and is therefore taken to index processing at the level of orthographic form. Family size affects the breadth of meaning-level connections activated when a morphological base is encountered, and is therefore taken to index processing at the level of semantic representation \citep{schreuderHowComplexSimplex1997, taftMorphologicalDecompositionReverse2004, taftRecognitionAffixedWords1979, dejongMorphologicalFamilySize2000, bertramEffectsFamilySize2000}.

Current accounts of morphological processing are designed to capture this pattern through different architectures — whether complex words are obligatorily decomposed into constituent morphemes, accessed via parallel whole-word and decompositional routes, or processed through learned form-to-meaning mappings without discrete morpheme representations at all \citep{schreuderModelingMorphologicalProcessing1995, dejongMorphologicalFamilySize2000, taft_lexical_1975, taftMorphologicalDecompositionReverse2004, rastleMorphologicalDecompositionBased2008, longtinMorphologicalDecompositionEarly2005, morrisSemanticTransparencyMasked2007, chuangDiscriminativeLearningLexicon}. Despite this architectural diversity, these accounts converge on a shared prediction: base frequency effects should arise earlier in processing than family size effects, reflecting a general precedence of form-level access over semantic integration. We return to how each account explains this precedence, and to what our results imply for adjudicating between them, in the General Discussion; here, the point is simply that this precedence is predicted across otherwise divergent architectures, which is what motivates using it to test the timing of OS and LP effects.

\subsection{Orthographic Sensitivity and Language Proficiency as Dissociable Dimensions}

Traditional models of word recognition typically characterise processing at the level of the skilled reader in general, abstracting over individual differences in favour of a uniform cognitive architecture. Yet readers do not differ along a single dimension of lexical skill. Andrews and colleagues tested readers' vocabulary knowledge and spelling ability and examined how these skills predicted morphological priming \citep{andrewsIndividualDifferencesSkilled2012, andrewsMorphologicalPrimingStronger2013}. They found that these two dimensions are dissociable and predict distinct patterns of morphological processing: readers with a semantic profile, defined by relatively higher vocabulary than spelling ability, show robust priming for transparent morphological pairs but little priming for opaque or form-only pairs, particularly at longer latencies, while readers with an orthographic profile, defined by relatively higher spelling than vocabulary, show sustained priming for opaque pairs that is at least as strong as for transparent pairs. \citet{andrewsMorphologicalPrimingStronger2013} interpret this as evidence that readers differ not merely in overall processing efficiency but in the relative weight they assign to form-based versus meaning-based information during lexical access.

If base frequency and family size index form-based and meaning-based aspects of lexical representation respectively, as established above, then readers who differ in their balance of orthographic and semantic processing should differ predictably in their sensitivity to each variable. To capture this variation continuously rather than as a categorical split, we derive two orthogonal dimensions, Language Proficiency and Orthographic Sensitivity, from a principal components analysis of vocabulary knowledge, spelling ability, and print exposure, following \citet{andrewsMorphologicalPrimingStronger2013}.

\subsection{The N250 and N400 as Indices of Form-Based and Meaning-Based Processing}

To track whether base frequency and family size effects emerge at distinct temporal phases, the present study uses two ERP components as dependent measures: the N250, an index of sublexical form processing, and the N400, an index of lexical-semantic integration. These components have been argued to index temporally distinct aspects of processing, from sublexical form analysis to semantic integration \citep{graingerWatchingWordGo2009}.

The N250 is a negative-going deflection peaking at approximately 250~ms post-stimulus onset. \citet{holcombTimeCourseVisual2006} propose that it reflects a phase of processing in which sublexical orthographic representations — that is, letters and letter clusters — are mapped onto higher-level form representations prior to full lexical access \citep{holcombTimeCourseVisual2006}. Evidence for this interpretation comes from its graded sensitivity to orthographic overlap between prime and target in masked priming: amplitude is most attenuated for full identity repetitions, intermediate for primes sharing all but one letter, and largest for unrelated pairs \citep{holcombTimeCourseVisual2006, holcomb_exploring_2007}. The N250 is also sensitive to sublexical phonological overlap, with \citet{graingerTimeCourseOrthographic2006} identifying separable early orthographic and later subcomponents, consistent with its role as an early, form-based phase of processing.

The N400 is a negative-going component peaking around 400~ms with a broad centro-parietal distribution \citep{kutasReadingSenselessSentences1980}. Its amplitude varies with semantic priming, cloze probability, and semantic richness, while remaining insensitive to physical and syntactic manipulations that leave meaning intact \citep{kutasThirtyYearsCounting2011}.

The N250 and N400 are dissociable not only by topography — frontal versus posterior \citep{holcombTimeCourseVisual2006} — but by a double dissociation across priming types: lexical status modulates the N400 but not the N250, since N400 priming is restricted to real words and does not extend to pseudowords \citep{kiyonagaMaskedCrossmodalRepetition2007}, while the reverse holds for semantic relatedness, which modulates the N400 but produces no N250 effect \citep{graingerWatchingWordGo2009}. This pattern confirms that the two components index qualitatively distinct processing phases, ruling out the interpretation of the N250 as merely an early onset of the N400.


	\subsection{Aims and Predictions}

	The present study has two complementary aims. The first is to establish whether base frequency and morphological family size exert their effects at distinct temporal stages of visual word recognition. The second, and primary, aim is to examine whether Orthographic Sensitivity and Language Proficiency modulate these effects selectively, at the processing stages where each variable is predicted to operate.

	The core prediction is a pattern of component-specific modulation: Orthographic Sensitivity should modulate base frequency effects at the N250, while Language Proficiency should modulate family size effects at the N400. If this pattern is observed, it would constitute direct electrophysiological evidence that the distinction between form-based and meaning-based morphological processing is a real axis of variation in the reading system, not merely a property of competing computational models. If instead both variables are modulated by the same individual difference dimension at the same stage, this would suggest that the two effects are less dissociable than current multi-stage accounts propose.

	The first analysis examines this prediction in real word recognition. Because word stimuli were not matched across morphological variables, base frequency, family size, and morphological complexity are treated as continuous predictors in a regression framework for reaction times; for ERP data, items were median-split into high and low groups separately for base frequency and family size. The two measures therefore carry somewhat different inferential weight: the RT regression speaks to the independent continuous contributions of each variable, while the ERP analyses speak to the direction and timing of effects at the extremes of each variable's distribution.

	The second analysis examines whether these modulations generalise to novel morphological structure, using morphologically complex nonwords. Morphologically structured nonwords (e.g., \textit{gasful}: real stem + real affix) were compared against matched controls (e.g., \textit{gasfil}: real stem + pseudoaffix), holding stem identity constant while manipulating suffix legitimacy. Because these items have no stored whole-word representations, they place maximal pressure on the system's generalisation mechanism and allow stem-level variables to be examined independently of whole-word familiarity. Orthographic Sensitivity is predicted to modulate early N250 responses to the suffix contrast, consistent with affix-led parsing: readers with more precise representations of real suffix forms should show greater N250 sensitivity to the distinction between real and pseudo-suffixes. Language Proficiency, building on the stem-led parsing possibility developed above, is predicted to modulate N250 responses through a distinct mechanism, and to predict later N400 sensitivity to whether a nonword's stem belongs to a large or productive morphological family, since it is this family membership that determines the richness of the semantic activation the form can support.
